\section{Conclusion and future work}\label{sec:final}

In this paper, we extend the paradigm of relation-changing logics to probabilistic relational structures. While traditional relation-changing logics address dynamic phenomena over discrete models, we adapt this framework to settings where transitions are governed by probabilities. Specifically, we introduce three novel relation-changing operators tailored to the probabilistic case and analyze their impact on the underlying semantics, namely, local sabotage, bridge, and swap modalities. These operators are inspired by~\cite{ArecesFH15} and are designed to represent transition deletion, addition, and reversal, respectively. To support reasoning in this setting, we also investigate appropriate notions of probabilistic bisimulation, showing how these operators interact with key behavioral equivalences. This work thus opens a new direction in the study of relation-changing logics, complementing existing proposals in paraconsistent, fuzzy, and graded contexts~\cite{CostaFM23,FFM25}.

This work is closely related to graded relation updates from~\cite{FFM25}, where classical relation-changing operators are extended to graded (fuzzy) relational structures. While that approach captures degree-based updates of accessibility, our proposal focuses on distributional updates over probabilistic transition systems, thus providing a complementary perspective on extending relation-changing logics beyond the classical setting. Moreover, unlike the approach in~\cite{FFM25}, we focus on a \emph{crisp} semantics, i.e., the truth value of a formula is ultimately true or false, whereas in the aforementioned work, formulas are interpreted with different degrees of truth.

We introduced the syntax, semantics, and examples for the three kinds of updates, and then studied the associated notions of bisimulation. We prove that these notions properly capture structural properties of the respective logics. Concretely, we establish invariance theorems and characterize appropriate Hennessy-Milner classes of models. We also present a case study to illustrate how our framework can be used for reasoning in real-life scenarios.  
Our developments can be considered a starting point for this perspective on probabilistic updates and their applications.

As future work, we aim to develop model-checking tools for probabilistic relation-changing logics. Although developed in a different context, rPRISM~\cite{FRMC19} -- a tool integrated with PRISM -- could serve as a starting point. We also plan to investigate the decidability status and complexity of the model-checking problem and develop an automatic model checker that integrates probabilistic sabotage, bridge, and swap operators. Furthermore, other reasoning tasks for probabilistic relation-changing logics remain unexplored. In particular, defining axiomatic systems, especially tableau-based methods inspired by, e.g.,~\cite{ArecesFH13,Fervari14,KozhemiachenkoS26}, and developing corresponding implementations extending, e.g.,~\cite{HoffmannA09,ArecesFHM16}, are promising directions for further research.

Another direction concerns the exploration of alternative semantics for rel\-ation-changing operators. In our approach, the operators are defined so that, after the model is changed, we move to a new accessible state (with an appropriate probability). An alternative would be to update the model without moving from the current state (see, e.g.,~\cite{GierasimczukKV09}). We also plan to investigate global versions of these operators, in which changes are not made locally but can occur anywhere in the transition probability matrix of the model. A further possibility is to define general probabilistic updates, as done in~\cite{ArecesFH15} for the non-probabilistic cases. Finally, restricted variants of the operators, designed for particular purposes as in epistemic logic, deserve further attention, particularly regarding their logical consequences for decidability and expressivity.

% Finally, we intend to explore the existence of a Hennessy-Milner theorem for finite models concerning our definition of bisimulation in the probabilistic setting, as well as in the model changing case.

